\begin{quote}\scriptsize\begin{verbatim}
TASK. You are given a legal issue and a passage from
  a United States federal
appellate opinion. Decide whether the opinion
  resolved that issue.

LABELS.
  DECIDED     The authoring court resolved the
    issue. It stated a holding, a
              conclusion, or a
                determination on the
                issue, whether or not
                that
              resolution was necessary
                to the judgment, and
                whether or not it
              was later characterized
                as dictum.
  UNRESOLVED  The authoring court did not resolve
    the issue for itself.

If and only if the label is UNRESOLVED, assign one
  sublabel.
  AVOIDED     The court declined to reach the
    issue because some other ground
              disposed of the case or
                the claim. Typical
                form: "because we
              affirm on qualified
                immunity grounds, we
                need not decide
                whether
              the search was lawful."
  ASSUMED     The court proceeded on a stated
    assumption about the issue and
              resolved the case on
                other grounds, so that
                the assumption did
              no independent work.
                Typical form:
                "assuming without
                deciding
              that the statute reaches
                this conduct, the
                evidence is
              insufficient."
  RESERVED    The court expressly held the issue
    open for future resolution,
              signalling that it
                remains available.
                Typical form: "we
                leave for
              another day whether
                Chevron survives in
                this context."

DECISION RULES.
  L1-1  ATTRIBUTION. The label describes what THIS
    opinion did. Language about
        what a different court did, whether
          a lower court, a sister circuit,
          or
        the Supreme Court, does not by
          itself make this opinion's
          disposition
        UNRESOLVED. "The district court
          concluded it need not decide X"
          tells
        you about the district court. Read
          on to see what this court did.
  L1-2  UNIT. In a cluster containing a majority
    opinion and separate opinions,
        each opinion is its own unit. A
          dissent that would decide an issue
          the
        majority avoided is DECIDED for the
          dissent and UNRESOLVED for the
        majority.
  L1-3  LATER RESOLUTION CONTROLS. A court may
    write "we need not decide X" and
        then decide X anyway, in the same
          opinion, in an alternative
          holding, in
        a footnote, or in a later section.
          Where the opinion does resolve the
        issue somewhere, the label is
          DECIDED. A court's
          characterization of
        what it need not do controls only
          where no independent resolution
        follows.
  L1-4  ASSUMPTION THAT DOES WORK. If the court
    assumes a proposition and the
        assumption is then treated as
          settled for the remainder of the
          analysis
        in a way that determines the
          outcome, the label remains
          UNRESOLVED and
        the sublabel is ASSUMED. Assuming a
          proposition is not deciding it,
          even
        when the assumption drives the
          result.
  L1-5  NO TRIGGER REQUIRED. UNRESOLVED does not
    require any particular phrase.
        A court can leave an issue open in
          ordinary prose.
  L1-6  PARTIAL RESOLUTION. The unit is the
    proposition named in the issue
        statement, not the broader legal
          question it belongs to. If the
          court
        resolves the anchored proposition
          but leaves a neighbouring one
          open,
        the label is DECIDED. If the court
          resolves a neighbouring
          proposition
        but leaves the anchored one open,
          the label is UNRESOLVED.
  L1-7  REMANDS. An issue remanded for the
    district court to address in the
        first instance is UNRESOLVED,
          sublabel AVOIDED, unless the
          appellate
        court also states the governing rule
          and applies it.
  L1-8  ABSTAIN. If the passage is too corrupt,
    too truncated, or too ambiguous
        to support either label, output
          UNCLEAR rather than guessing.
  L1-9  ISSUES THIS OPINION NEVER TAKES UP.
    Sometimes the only reason an issue
        is in front of you is that the
          opinion mentioned, in a
          parenthetical or
        in a description of some other
          court's work, that a different
          court did
        not decide it, and this opinion
          never reaches the issue in its own
        right. That is neither DECIDED nor
          UNRESOLVED for this opinion.
          Output
        UNCLEAR. This rule does not apply
          where the opinion goes on to
          engage
        with the issue itself: there, L1-1
          and L1-3 govern.
\end{verbatim}\end{quote}

\noindent Version 1.1. The single permitted revision after the pilot, and the layer-two and issue-statement guidelines, are in the released package.
